\section{A hybrid renormalization procedure}
\label{sec:hybrid}

As explained in the previous section,
the LaMET expansion starts from calculating the coordinate-space
correlation functions $\tilde h(z,P^z)$ at large momentum $P^z$ and for the whole range
of distance $-\infty<z<\infty$.
On a discrete
lattice with spacing $a$, the nonlocal quark bilinear operator that defines $\tilde h(z,P^z)$ in \eq{esme} can be multiplicatively renormalized as~\cite{Ji:2017oey,Ishikawa:2017faj,Green:2017xeu}
\begin{align}\label{eq:ren}
&\left[\bar{\psi}(z)\Gamma W(z,0)\psi(0)\right]_{\rm B} \nn\\
&= e^{\delta m|z|}Z(a) \left[\bar{\psi}(z)\Gamma W(z,0)\psi(0)\right]_{\rm R}\,,
\end{align}
up to lattice artifacts~\cite{Constantinou:2017sej,Chen:2017mie}. Here the operator on the l.h.s. is defined in terms of bare fields and couplings, denoted by the subscript ``B'', while the operator on the r.h.s. is renormalized and denoted by the subscript ``R''. Without an explicit
statement, we always assume that the renormalized correlations are eventually defined in the $\MS$ scheme before a Fourier transformation is made to the momentum space.
There are both $z$-independent
logarithmic and $z$-dependent linear divergences.
The former arises from the renormalization
of quark and gluon fields as well as the vertices at the endpoints of the Wilson line, which is included in the factor $Z(a)$, while
the latter comes from the Wilson-line self-energy, which is factored into the exponential $e^{\delta m|z|}$ with $\delta m$ being the ``mass correction'' .

A number of
proposals have been made in the literature~\cite{Ishikawa:2016znu,Chen:2016fxx,Monahan:2016bvm,Radyushkin:2017cyf,Constantinou:2017sej,Green:2017xeu,Stewart:2017tvs,Braun:2018brg,Li:2020xml} to renormalize the above lattice correlation functions $\tilde h(z, a, P^z)$, among which the RI/MOM scheme has frequently
been used~\cite{Constantinou:2017sej,Stewart:2017tvs}. In this approach,
one calculates the matrix elements (amputated Green's function) $Z(z,-p^2,a)$ of the bilocal operators $O(z,a)$ in a deep Euclidean state with momentum squared $-p^2\gg\Lambda_{\rm QCD}^2$ in a
fixed gauge, and then defines $\overline{\rm MS}$ operators
as,
\begin{align}
             O_{\overline{\rm MS}}(z,\mu) \equiv Z_{\overline{\rm MS}}(z,-p^2,\mu)\frac{O(z,a)}{Z(z,-p^2,a)},
\end{align}
where $Z_{\overline{\rm MS}}$ converts the RI/MOM renormalized result to the $\overline{\rm MS}$ scheme. The gauge and $-p^2$ dependences cancel between two $Z$-factors.
The r.h.s. has a proper continuum limit $a\to 0$ without divergences.

However, while the RI/MOM approach is justified
for local operators, it has potential problems when
applied to nonlocal ones.
For instance, when $z$ becomes large,
$Z_{\overline{\rm MS}}(z,-p^2,\mu^2)$ contains IR
logarithms of $z$ and the perturbative calculation
of $z$-dependence is not reliable. Moreover,
although the RI/MOM factor $Z(z,-p^2,a)$ helps to cancel
the lattice UV divergences, the composite operator
at large-$z$ contains non-perturbative physics as well.
Therefore, both $Z$-factors contain non-cancelling
non-perturbative effects which alter the IR properties
of $O(z)$. Thus, the RI/MOM renormalization scheme is not reliable
at large-$z$. Moreover, when gluon distributions
are involved, it requires external off-shell gluon states which
bring in potential mixing with gauge-variant operators and make things much more complicated~\cite{Wang:2019tgg}.

In addition to the renormalization issues at large distances, there are also
subtleties for renormalization at short distances. While the standard renormalization
of a bilocal operator makes it finite at any non-vanishing
$z$, it becomes divergent in the $z\to 0$ limit. On the other hand, if one performs a resummation of the large
logarithms at small $z$, the result vanishes at $z=0$. However, the lattice
result at $z=0$ approaches to the matrix element of the vector current.
Clearly, the two limits, $a\to 0$ and $z\to 0$,
are not interchangeable.

To resolve these issues, we propose in this section a
hybrid scheme to renormalize the correlation functions. The key point of this scheme is that we separate the correlations
at short and long distances and renormalize them separately, and match both procedures
at an intermediate distance $z_{\rm S}$. The
matching point must lie within $[0,z_{\rm LT}]$ where
the leading-twist (LT) approximation for the correlation
operator is valid. Discussions on the value of
$z_{\rm LT}$ can be found in Sec. V.

\subsection{Renormalization at short distance $0\le |z|\le z_{\rm S}$}
\label{sec:short}

To renormalize $\tilde h(z,a, P^z)$ for $0\le |z|\le z_{\rm S}$, particular attention shall be paid to the behavior of the correlation functions in the limit $z\to 0$.

In the continuum $\MS$ sheme, the $z\rightarrow 0$ limit is not smooth and
additional logarithmic UV divergences $\sim \ln z^2$ arise which
when resummed yield zero.  However, this is not the case for the lattice matrix element
$\tilde h(z,a, P^z)$. For finite lattice spacing and non-vanishing $z$,
$\tilde h(z,a, P^z)$ includes UV divergences related to the wave function
renormalization of the bare fields, of the form
$\alpha_s(a)\ln (z^2/a^2)$.
At small $z$, particularly when $z=0$ or $a$, $\tilde h(z,a, P^z)$ has discretization effects
and is related to the
lattice-regulated local matrix element ${\bar\psi}\Gamma \psi$.
In particular, when $\Gamma=\gamma^\mu$,
${\bar\psi}\gamma^\mu \psi$ is conserved and its matrix element is finite
in the $a\to 0$ limit. A function demonstrating this interesting
interplay between lattice regulator and small physical distance
is $\ln [(z^2+a^2)/a^2]$.

The above discrepancy in the small-$z$ regime can be removed through a perturbative conversion between lattice regularization and the continuum $\MS$ scheme, which, however, is known to converge slowly. Instead, a more efficient strategy is to cancel the $\ln z^2$-dependences through lattice renormalization, which corresponds to a scheme ``X'' that is different from $\MS$.
As long as $z$ is in the leading-twist region $|z|\le z_{\rm S}$ where $z_{\rm S}$ is
smaller than $z_{\rm LT}$, the difference between the X-scheme and $\MS$ can be calculated in perturbation theory.

For example, the X-scheme can be implemented by forming the ratio of $\tilde h(z,a, P^z)$ and another
matrix element of the same operator $O(z,a)$,
\begin{align}\label{eq:ratio}
 \frac{\tilde h(z,a,P^z)}{Z_X(z,a)}\,,~~~~~~{\rm for\ }|z|\le z_{S}\,,
\end{align}
where the renormalization factor $Z_X$ corresponds to different choices of the matrix element.
Possible choices for $Z_X$ include
\begin{itemize}
\item{Amputated Green's function of $O(z,a)$ in a single-particle deep Euclidean state, fixed in a particular gauge, e.g., Landau gauge, which defines the RI/MOM-type of schemes~\cite{Constantinou:2017sej,Stewart:2017tvs,Alexandrou:2017huk,Chen:2017mzz}.}
\item{Matrix element of $O(z,a)$ in a hadron state with $P^z=0$, depending on applications~\cite{Radyushkin:2017cyf,Orginos:2017kos}.}
\item{Vacuum ($|\Omega\rangle$) expectation value of $O(z,a)$~\cite{Braun:2018brg,Li:2020xml}.}
\end{itemize}
In the second and third option, the matrix elements are gauge invariant,
and therefore no gauge fixing is needed. For the third option, the quantum numbers of the operator must be the same as those of the vacuum.
As discussed above, $z_{\rm S}$ has to be smaller than $z_{\rm LT}$, which is estimated in Sec. V to be about
$0.25 \sim 0.33$~fm. Of course, the stability of the final result with respect to small variations of $z_{\rm S}$ shall be explicitly verified.

Due to the multiplicative renormalizability of the operator $O(z,a)$, all UV divergences cancel in the ratio in \eq{ratio}, thus allowing us to take the continuum limit,
\begin{align}\label{eq:ratio2}
\lim_{a\to0} \frac{{\tilde h}(z,a,P^z)}{Z_X(z,a)} = \frac{{\tilde h}(z,\epsilon,P^z)}{Z_X(z,\epsilon)} \equiv\tilde h^X(z,P^z)\,,
\end{align}
where the term after the first equal sign refers to a $\MS$
calculation of the same ratio with $\epsilon$ corresponding to dimensional regularization $d=4-2\epsilon$ in the continuum theory. In the limit $|z|\to0$, the $\ln z^2$-dependence is independent of the external state, so it cancels in the ratio, making the latter finite at $z=0$.
Moreover, in the leading-twist region $z\le z_{\rm S} \le z_{\rm LT}$, we can perturabtively match the ratio for any X-scheme to the LF correlation $h(\lambda,\mu)$ through the coordinate-space factorization formula~\cite{Izubuchi:2018srq,Radyushkin:2017lvu}
\begin{align}\label{eq:coordfact}
\tilde h^X(\lambda, P^z) & =\int_{0}^1 d\alpha\ {\cal C}^X\Big(\alpha,\frac{\lambda^2\mu^2}{(P^z)^2}\Big) h(\alpha\lambda,\mu) \nn\\
&\qquad + {\cal O}(z^2\Lambda_{\rm QCD}^2)\ ,
\end{align}
where ${\cal C}^X$ is the matching coefficient, and we have suppressed its dependence on the renormalization scale in the X-scheme such as the RI/MOM. Also, higher-twist contributions have been suppressed in the above equation.

The one-loop matching coefficient for the second option for $Z_X$ has been obtained in Refs.~\cite{Zhang:2018ggy,Radyushkin:2017lvu,Izubuchi:2018srq}, and that for the RI/MOM scheme can be extracted from Ref.~\cite{Constantinou:2017sej,Izubuchi:2018srq}. The two-loop results for both the second and third option have been calculated as a series expansion in Ref.~\cite{Li:2020xml}.

\eq{coordfact} also has an equivalent form in momentum space~\cite{Izubuchi:2018srq}, with the matching coefficient calculated at one-loop~\cite{Stewart:2017tvs,Liu:2018uuj,Liu:2018hxv} and two-loop~\cite{Chen:2020ody} orders.
The two-loop matching coefficient for $Z_X$ defined by $P^z=0$ matrix element can be extracted from Refs.~\cite{Chen:2020ody,Li:2020xml}.


\subsection{Renormalization at large distances $z\ge z_{\rm S}$}
\label{sec:long}

At large distance $z\ge z_{\rm S}$, UV renormalization
needs a careful assessment because both the RI/MOM and the ratio scheme will introduce
undesired non-perturbative effects.
The only renormalization approach that will not introduce such extra non-perturbative physics
is the explicit and separate subtraction of linear divergences (or $\delta m$)
and logarithmic divergences~\cite{Chen:2016fxx,Ji:2017oey,Ishikawa:2017faj,Green:2017xeu},
which in principle can be done using the auxiliary field method~\cite{Green:2017xeu,Green:2020xco}.

To calculate the mass renormalization $\delta m$
of the Wilson line, there exist many
suggestions in the literature. Here we provide probably an incomplete list:

\begin{itemize}
\item{One can fit the hadron matrix element
at large $z$, where the dominant decay is
\begin{equation}\label{eq:mass1}
  \tilde h(z,a, P^z) \sim \exp(-\delta m |z|) \ .
\end{equation}
$\delta m$ can be obtained by fitting the ratio $\ln (\tilde h(z+a,a, P^z)/\tilde h(z,a, P^z))$
to a constant in $z$ at large $z$. Of course, the result has to be independent of $P^z$,
e.g., one can choose $P^z=0$. Alternatively, one can fit $\ln \tilde h(z,a, P^z)$ to the $1/a$
dependence all $z$. This method has yet to be studied using real lattice data.}
\item{One can use the single-quark Green's function
as in the RI/MOM renormalization factor,
\begin{align}
          Z(z,-p^2,a)&=\int d^4x d^4y\ e^{ip\cdot (x-y)}\nn\\
          &\times  \langle \Omega|T\psi(x) O(z,a) \bar{\psi}(y)|\Omega\rangle\,,
\end{align}
which asymptotically goes like $Z(z) \sim \exp(-\delta m |z|)$. This matrix element needs a fixed gauge,
and has been studied in Refs.~\cite{Izubuchi:2019lyk,Huo:2019vdl}.}
\item{On can also use the vacuum matrix element of $O(z,a)$
\begin{equation}
        S(z) = \langle \Omega|O(z,a)|\Omega\rangle
\end{equation}
which is gauge invariant. $S(z)$ again at very large $z$ behaves like
$S(z) \sim \exp(-\delta m |z|)$. This has been considered in Refs.~\cite{Braun:2018brg,Li:2020xml}.
}
\item{Also the gauge-invariant Polyakov loop leads to the static potential between two heavy quarks. There exists a large number of references on this approach, see, e.g.,~\cite{Appelquist:1977tw,Schroder:1999xf,Smirnov:2009fh,Philipsen:2002az,Jahn:2004qr}.}
\item{One can also calculate the vacuum expectation value
of the Wilson line $W(z)$ directly in a fixed gauge,
and again $\langle W(z) \rangle \sim \exp(-\delta m |z|)$ at large distance. This has been considered in Refs.~\cite{Green:2017xeu,Green:2020xco} using the auxiliary field method.}
\end{itemize}
It is worth pointing out that, although all the proposals above work in principle, different practical issues may arise in their lattice realization such that some may work better than the others.

The mass renormalization $\delta m$ is gauge-independent, just like the pole mass
of a quark~\cite{Kronfeld:1998di}. In the above suggestions where no gauge fixing is
needed, this is obviously true. In the cases where a gauge-fixing is needed,
one can demonstrate that the results in any other gauge are the same by
constructing appropriate gauge-invariant operators~\cite{Philipsen:2001ip}.
Despite being gauge-independent, $\delta m$ will depend on the specific action used in Monte Carlo simulations
and on the definition of the matrix elements above.
In the cases of vacuum matrix elements, $\delta m$ may be interpreted
as the non-perturbative pole mass in certain gauges~\cite{Philipsen:2001ip}.

The $\delta m$ calculated from all the matrix elements above will have the following dependence
on the lattice spacing $a$,
\begin{equation}
   \delta m = m_{-1}(a)/a + m_0 \ ,
\end{equation}
where $m_{-1}$(a) is the coefficient of the power divergence, which is independent of the specific matrix element. The $a$-independent term $m_0$
has a more complicated origin. It can arise
from various sources:
\begin{itemize}
\item{Renormalon effect: In principle, $m_{-1}(a)$ can be calculated perturbatively, and is $2\pi\alpha_s(a)/3$ at leading
order, similar to the mass counterterm in the Wilson
formulation of fermions. However, the perturbation series is not convergent. When truncated at order $n\sim 1/\alpha_s(a)$, the
perturbation series has an uncertainty of order $a\Lambda_{\rm QCD}$, which generates a contribution to $m_0$~\cite{Ji:1995tm,Beneke:1998ui,Bauer:2011ws,Bali:2013pla}.
This means that in non-perturbative fitting, $m_{-1}(a)$ is determined only
with an uncertainty of the order of $a\Lambda_{\rm QCD}$. Thus, an additional
$m_0$ contribution of order $\Lambda_{\rm QCD}$ is expected.}
\item{Pole mass: For certain matrix elements, like vacuum elements of bilocal operators, the $z$
dependence can be viewed as originating from the pole mass of a meson consisting of an infinitely-heavy quark
and a light one. In this case, $\delta m$ is the pole mass apart from the linear divergence.}
\item{Finite $P^z$ effects: The correlation function at finite $P^z$ has a long-range correlation,
$\exp(-\lambda/\xi(P^z))$, where $\xi(P^z)$ is the correlation length. This contribution is included in $m_0$.}
\item{Fitting effect. Since the data is always in finite $z$ and $a$ where the exponential decay
cannot always be separated from an algebraic decay, there are fitting uncertainties
contributing to $m_0$ as well as the separation between $m_0$ and $m_{-1}$.}
\end{itemize}
To summarize, $m_0$ depends on the lattice matrix-element used and the fitting procedure~\cite{Zhang:2017bzy,Chen:2017gck,Izubuchi:2019lyk,Huo:2019vdl}.

In Fig.~\ref{fig:deltam} we show, as an example, the values of $m_{-1}$ and $m_0$ determined from the quark RI/MOM renormalization factor calculated at the scale $\mu_R=1.8~ {\rm GeV}$ and $p^z_R=0$, using the four ensembles with $a\approx \{0.045, 0.06, 0.09, 0.12\} {\rm fm}$ and $310$ MeV pion mass from MILC collaboration~\cite{Bazavov:2012xda}. Inspired by the asymptotic behavior at large $z$ to be studied in Sec.~\ref{sec:regge}, we use the following simplified form
\begin{equation}
e^{-(\frac{m_{-1}}{a}+m_0)|z|}\frac{c_1}{|z|^{d_1}}
\end{equation}
to fit the renormalization factors at four different lattice spacings. It is worth pointing out that $m_{-1}$ starts from $O(\alpha_s)$ and we therefore also include the dependence of the coupling on $a$ in the fitting. For $a\approx 0.12$ fm, the fitted results for the coefficients $m_{-1}$ and $m_0$ are
\begin{equation}
m_{-1}=0.234\pm 0.012, \ \ m_0=(350\pm 60)\, {\rm MeV}.
\end{equation}

\begin{figure}[tbp]
\centering
\includegraphics[width=0.9\columnwidth]{images/npr_deltam_new}
\vspace*{-2em}
\caption{Fitting of the quark RI/MOM renormalization factor calculated using four ensembles with lattice spacings $a\approx \{0.045, 0.06, 0.09, 0.12\} {\rm fm}$ and $310$ MeV pion mass from MILC collaboration~\cite{Bazavov:2012xda}. }
\label{fig:deltam}
\end{figure}

In principle, one can choose to subtract the power divergent piece only, namely $m_{-1}/a$.
The less-well-determined $m_0$ term can be left in the lattice matrix elements, the
momentum expansion will take care of the rest. Indeed, the difference between subtracting
different $m_0$
is $O(1/P^z)$ effect, as has been demonstrated
perturbatively in~\cite{Xiong:2013bka}.
More precisely, assuming there are two quasi-LF correlations that define the quasi-PDFs and differ from each other by a factor  $e^{-m|z|}$ with $m \sim \Lambda_{\rm QCD}$,
\begin{equation}
\tilde h_1(z,P^z)=\tilde h_2(z,P^z)e^{-m|z|} \ ,
\end{equation}
then after Fourier transforming into momentum space, they are related by
\begin{align}
&\tilde f_1(y,P^z)-\tilde f_2(y,P^z)\nonumber \\
&=\frac{1}{\pi}\int_{-\infty}^{\infty} dy' \frac{\delta}{\delta^2+(y-y')^2} \left[\tilde f_2(y',P^z)-\tilde f_2(y,P^z) \right]\ ,
\end{align}
with $\delta=m/P^z$. If the $\tilde f$'s are square integrable and their first order derivatives are continuous, one can show that as $\delta \rightarrow 0$,
\begin{align}
|\tilde f_1(y,P^z)-\tilde f_2(y,P^z)|\sim \delta \ .
\end{align}
Therefore, the ambiguity of different schemes disappears in the large $P^z$ limit.

However, this is still unsatisfactory because it appears that, due to non-perturbative
effects from the linear divergence, the LaMET expansion will be
an expansion in powers of $M/P^z$ instead of $(M/P^z)^2$, which
will significantly reduce the speed of convergence. Here
we consider possible ways to overcome this deficiency. Recall that the LaMET expansion in $(M/P^z)^2$ is
made in the $\MS$ scheme where no linear divergence exists, in a general scheme this expansion might
contain odd powers in $1/P^z$. Therefore, there is a way to choose $m_0$
such that the condition of the $\overline{\rm  MS}$ scheme
\begin{equation}
      \delta m_{\overline{\rm MS}} = 0
\end{equation}
is met with non-perturbative calculations.
We shall denote such a value of the subtracted mass
by \begin{equation}
\delta m_c = m_{-1}/a + m_{0c}\ .
\end{equation}
where $m_{0c}$ can be determined by matching the matrix element $\tilde h_1(z, P^z, a)$ on lattice
to the $\overline{\rm MS}$ result when $z$ is small (around $z_S$) and
QCD perturbative theory works. An alternative strategy
is to vary $m_{0}$ in a certain range near $\Lambda_{\rm QCD}$,
and identify the value $m_{0c}$ for which the linear term in $1/P^z$ in $\tilde f(y, P^z)$ vanishes.
This is like searching for the critical value of $\kappa_c$ for Wilson
fermions for which a similar power divergent
bare quark mass appears~\cite{Lepage:1992xa}.

The mass-subtracted operator $O(z,a)e^{-\delta m(a)|z|}$
has no power divergence, but still has logarithmic
dependence on $a$.
These remnant logarithmic divergences are independent of $z$ and can be renormalized, in principle, using the auxiliary field method~\cite{Green:2017xeu,Green:2020xco}.
However, a more convenient option in practice is to fix the renormalization constant $Z_{\rm hybrid}(a)$ by directly
matching the renormalized matrix elements of $O(z,a)$ at $z=z_{\rm S}$ from the short and long distance regimes, which is essentially a continuity condition,
\begin{align}
Z_{\rm hybrid} e^{-\delta m |z_{\rm S}|}\langle P| O(z_{\rm S},a) | P \rangle  = \frac{\langle P| O(z_{\rm S},a) | P \rangle }{Z_X(z_{\rm S},a) }\,,
\end{align}
which leads to
\begin{align}
Z_{\rm hybrid}(z_{\rm S},a) = e^{\delta m |z_{\rm S}|}/{Z_X(z_{\rm S},a)  }\,.
\end{align}
In this way, one only has to calculate $\delta m$. Of course, one needs to vary $z_{\rm S}$ to check whether the final result is stable.

The matching coefficient ${\cal C}_{\rm hybrid}$ for the long distance regime is related to that for the X-scheme. For example, if one adopts the $P^z=0$ matrix element for renormalization~\cite{Zhang:2018ggy,Radyushkin:2017lvu,Izubuchi:2018srq}, then
\begin{align}\label{eq:hybridm1}
&{\cal C}_{\rm hybrid}(\alpha, z^2\mu^2, z^2/z_{\rm S}^2) = {\cal C}_{\rm ratio}(\alpha, z^2\mu^2) \nn\\
&\qquad + \delta(1-\alpha){\alpha_sC_F\over 2\pi}{3\over 2} \ln {z^2\over z_{\rm S}^2}\ \theta(|z|-z_{\rm S})\,.
\end{align}
 However, due to the logarithms of $z^2\mu^2$ and $z^2/z_{\rm S}^2$, the above matching coefficient is only valid for $|z| \ll \Lambda_{\rm QCD}^{-1}$, otherwise one has to resum the large logarithms for $|z|\sim \Lambda_{\rm QCD}^{-1}$ by evolving $\alpha_s$ to a highly nonperturabtive regime. Since our ultimate goal is to Fourier transform the final result to obtain the PDF, this will introduce uncontrolled sytematics.

To have a clearer way of separating the perturbative and non-perturbative regimes, we can perform the matching in momentum space, where nothing prevents using the correlations at large $z$, provided that $yP^z$ is sufficiently large. In principle, we should first convert the hybrid scheme to the $\MS$ scheme---where the factorization formula was proven~\cite{Ma:2014jla,Ji:2020ect}---in coordinate space with the conversion factor
\begin{align}\label{eq:hybrid3}
    {\cal Z}(z,z_S,\mu) & = {Z_{\MS}\over Z_{\rm hybrid}} =\frac{\theta(z_{\rm S}\!-\!|z|)}{Z^{\overline{\rm MS}}_{X}(z,\mu)}
    +\frac{\theta(|z|\!-\! z_{\rm S})}{Z^{\overline{\rm MS}}_{X}(z_S,\mu)}\,,
\end{align}
where $1/Z^{\overline{\rm MS}}_{X}$ converts the ``$X$'' scheme to $\MS$, and for \eq{hybridm1}
\begin{align}\label{eq:zmsr}
    Z^{\overline{\rm MS}}_{\rm ratio} = 1 - {\alpha_sC_F\over 2\pi}\left[{3\over 2}\ln{z^2\mu^2\over 4e^{-2\gamma_E}}+ {5\over2}\right]\,.
\end{align}
The conversion factor ${\cal Z}$ is perturbative for all $z$ as it does not include $\ln(z^2)$ at large distance $|z|>z_{\rm S}$. Then we can Fourier transform the $\MS$ quasi-LF correlation and match it to the PDF in momentum space.

Since the scheme conversion is perturbative for all $z$, we can also do the Fourier transform first and directly match the hybrid scheme quasi-PDF to the PDF, and the matching coefficient $C_{\rm hybrid}$ is given by the double Fourier transform from \eq{hybridm1},
\begin{align}\label{eq:hybridm2}
&C_{\rm hybrid}(\xi, \mu^2/(p^z)^2, z_{\rm S}^2\mu^2) = C_{\rm ratio}(\xi, \mu^2/(p^z)^2) \nn\\
&\quad+{\alpha_sC_F\over 2\pi}{3\over 2} \left[-{1\over |1-\xi|_+}+ \frac{2 \text{Si}((1-\xi) \lambda_{\rm S})}{\pi  (1-\xi)} \right]\,,
\end{align}
where $C_{\rm ratio}$ can be found in~\cite{Izubuchi:2018srq}, $\xi=y/x$, and $\lambda_{\rm S}=z_{\rm S}p^z$ with $p^z=xP^z$ being the parton momentum. The plus function is defined as
\begin{align}
	{1\over |1-\xi|_+} &\equiv \lim_{\beta\to0^+}\left[ {\theta(|1-\xi|-\beta)\over |1-\xi|} + 2\delta(1-\xi)\ln\beta\right] \,.
\end{align}

We can also derive the corresponding scheme conversion factor and matching coefficient for using RI/MOM scheme in the short-distance renormalization. In the limit of $-z^2_{\rm S}p^2 \ll 1$, the RI/MOM renormalization factor $Z(z,-p^2)$ for $|z|<z_{\rm S}$ is equal to that of the zero-momentum matrix element, so the results are the same as those in \eqs{hybrid3}{hybridm2}. However, if $-z^2_{\rm S}p^2$ is finite, then one needs to use the results from Refs.~\cite{Constantinou:2017sej,Stewart:2017tvs} to derive the scheme conversion factor and matching coefficient.
Finally, since the result of the PDF must be independent of the lattice renormalization scheme, we can try different short-distance schemes and check if they are consistent with each other.

In momentum space, the matching coefficient includes the logarithm of $\mu/(yP^z)$ which becomes non-perturbative for $y\sim \Lambda_{\rm QCD}/P^z$. This is consistent with the power counting parameter $\Lambda^2_{\rm QCD}/(yP^z)^2$. Therefore, the nature of the systematic uncertainties is clear, and we can only improve precision at
small $x$ by pushing to higher $P^z$.
